\documentclass[10pt,twocolumn]{article}
\usepackage{amsmath,amssymb,booktabs,graphicx,hyperref,geometry}
\geometry{margin=0.85in,columnsep=0.25in}
\hypersetup{colorlinks=true,linkcolor=blue,citecolor=blue,urlcolor=blue}
\title{Quantum-Inspired vs.\ Classical Solvers: A Reproducible Benchmark}
\author{Anonymous}
\date{\today}
\begin{document}
\maketitle
\begin{abstract}
We present a reproducible benchmark comparing GQIR, SimCIM, SQA, SA, and CP-SAT on three QUBO problems across five sizes ($n=20$--$100$). With $3{,}180$ runs, we show SimCIM and SQA achieve near-zero gap on MaxCut, while classical SA dominates random QUBO ($0.12\%$ at $n=100$). CP-SAT achieves $0\%$ gap on MaxCut at all sizes. The QUBO penalty formulation for constrained problems is fundamentally ill-conditioned (gaps $>10^3\%$ regardless of penalty scaling). All quantum-inspired results are from classical simulation; no quantum hardware was used.
\end{abstract}

\section{Introduction}
This paper benchmarks quantum-inspired algorithms (SimCIM, SQA, GQIR) against classical SA and CP-SAT on QUBO-formulated optimization problems.

\section{Methods}
All problems minimize $x^\top Qx$ for $x\in\{0,1\}^n$.
\textbf{Random QUBO:} $Q_{ij}\sim\mathcal{U}[-1,1]$, density 0.3.
\textbf{MaxCut:} $Q_{ii}=-\frac12\sum_{j\neq i}w_{ij}$, $Q_{ij}=\frac12 w_{ij}$.
\textbf{Knapsack:} Penalty formulation $Q_{ii}=-v_i+P(w_i^2-2Cw_i)$, $Q_{ij}=2P w_i w_j$.

Algorithms: GQIR (gradient-based relaxation), SimCIM (phase dynamics), SQA (path-integral MC), SA (Metropolis), CP-SAT (OR-Tools, 30~s timeout for $n>20$).

\section{Results}
\begin{table*}[t]\centering\small
\caption{Mean gap (\%). CP-SAT$^\dagger$ = 30~s timeout.}
\label{tab:main}
\begin{tabular}{lcrrrr}
\toprule
Algorithm & Problem & $n$ & Gap(\%) & Succ(\%) & Time(s)\\
\midrule
SimCIM & MaxCut & 20 & $\approx$0 & 100 & 0.02\\
SimCIM & MaxCut & 100 & $\approx$0 & 100 & 0.03\\
SQA & MaxCut & 20 & 0.05 & 100 & 0.04\\
SQA & MaxCut & 100 & 0.02 & 100 & 0.16\\
SA & MaxCut & 20 & 0.64 & 74 & 0.01\\
SA & MaxCut & 100 & \textbf{0.03} & 100 & 0.03\\
CP-SAT$^\dagger$ & MaxCut & 20 & 0.00 & 100 & 2.81\\
CP-SAT$^\dagger$ & MaxCut & 100 & 0.00 & 100 & 30.05\\
GQIR & MaxCut & 20 & 3.39 & 0 & 0.04\\
GQIR & MaxCut & 100 & 6.90 & 0 & 0.03\\
\addlinespace
SA & RndQUBO & 20 & 2.23 & 80 & 0.01\\
SA & RndQUBO & 100 & \textbf{0.12} & \textbf{98} & 0.03\\
SQA & RndQUBO & 20 & 7.98 & 22 & 0.04\\
SQA & RndQUBO & 100 & 1.00 & 62 & 0.17\\
CP-SAT$^\dagger$ & RndQUBO & 20 & 4.91 & 70 & 0.03\\
CP-SAT$^\dagger$ & RndQUBO & 100 & 2.70 & 20 & 30.03\\
GQIR & RndQUBO & 20 & 23.03 & 30 & 0.04\\
GQIR & RndQUBO & 100 & 9.17 & 0 & 0.03\\
\bottomrule
\multicolumn{5}{l}{\footnotesize $^\dagger$ Feasible under 30~s timeout.}
\end{tabular}
\end{table*}

\subsection*{MaxCut}
SimCIM achieves $<10^{-14}\%$ gap. CP-SAT matches at $0\%$. SQA is $<0.1\%$. SA improves with $n$ ($0.64\%\to0.03\%$). GQIR shows $3$--$7\%$ gap.

\subsection*{Random QUBO}
Classical SA dominates ($0.12\%$ at $n=100$). CP-SAT is competitive ($0.79$--$4.91\%$) but encoding overhead hurts at $n=100$. SimCIM fails ($>130\%$). SQA improves with size ($8\%\to1\%$).

\subsection*{Penalty Pathology}
All heuristics produce gaps $>10^3\%$ on knapsack, regardless of penalty multiplier $\alpha$. This is a fundamental limitation of QUBO penalty encoding.

\section{Conclusion}
(1) SimCIM/SQA excel on MaxCut. (2) SA dominates unstructured QUBO. (3) CP-SAT is an excellent reference for MaxCut. (4) QUBO penalty encoding is fundamentally ill-conditioned.

\begin{thebibliography}{5}
\bibitem{kirkpatrick1983} S. Kirkpatrick et al., \emph{Science}, 1983.
\bibitem{leleu2021} T. Leleu et al., \emph{Phys. Rev. A}, 2021.
\bibitem{openjij} OpenJij Project, \url{https://openjij.github.io/}
\bibitem{ortools} Google OR-Tools, \url{https://developers.google.com/optimization/cp/cp_solver}
\bibitem{lucas2014} A. Lucas, \emph{Front. Phys.}, 2014.
\end{thebibliography}
\end{document}
